% chapters/ch05.tex —— 第 5 章：经典迭代法
% 源书：Chapter 5 Classical Iterative Methods，书页 60–75（物理页 74–89）
% 公式编号 (5.1)–(5.11)；定理 5.1–5.13；定义 5.1–5.5；引理 5.1；推论 5.1；例 5.1；习题 1–14
\chapter{经典迭代法}
\label{ch:iterative}

我们研究求解 $A\mathbf{x}=\mathbf{b}$ 的经典迭代法. 迭代法最初由 Gauss 于 1823 年, Liouville 于 1837 年和 Jacobi 于 1845 年提出, 它与高斯消去法等直接法有很大不同. 参见 \cite{src:2}.

若矩阵 $A$ 相当大, 则基于 $A$ 的 $LU$ 分解的直接法在计算时间与计算机存储两方面都变得难以承受. 在偏微分方程离散化等实际问题中, 矩阵的阶数可高达数十万. 对这类问题, 直接法变得不切实际. 此外, 大多数大型问题都是稀疏的, 而在 $LU$ 分解过程中稀疏性通常会丢失. 因此, 在 $LU$ 分解结束时我们不得不面对一个含有大量非零元素的大矩阵, 此时存储成为关键问题. 对这类问题, 我们可以使用一类称为迭代法的方法. 本章只讨论若干经典迭代法.

我们指出, 经典迭代法的缺点是收敛速度可能很慢, 甚至可能发散, 而且需要寻找合适的停止准则.

\section{Jacobi 方法与 Gauss-Seidel 方法}
\label{sec:5-1}

\subsection{Jacobi 方法}
\label{subsec:5-1-1}

考虑如下线性方程组
\[
  A\mathbf{x}=\mathbf{b},
\]
其中 $A=[a_{ij}]\in\R^{n\times n}$. 可以把矩阵 $A$ 写成如下形式
\[
  A=D-L-U,
\]
其中
\[
  D=\diag(a_{11},a_{22},\ldots,a_{nn}),
\]
\[
  L=
  \begin{bmatrix}
    0 & & & & \\
    -a_{21} & 0 & & & \\
    -a_{31} & -a_{32} & 0 & & \\
    \vdots & \vdots & \ddots & \ddots & \\
    -a_{n1} & -a_{n2} & \cdots & -a_{n,n-1} & 0
  \end{bmatrix},
\]
以及
\[
  U=
  \begin{bmatrix}
    0 & -a_{12} & -a_{13} & \cdots & -a_{1n}\\
      & 0 & -a_{23} & \cdots & -a_{2n}\\
      & & \ddots & \ddots & \vdots\\
      & & & 0 & -a_{n-1,n}\\
      & & & & 0
  \end{bmatrix}.
\]
于是容易看出
\[
  \mathbf{x}=B_J\mathbf{x}+\mathbf{g},
\]
其中
\[
  B_J=D^{-1}(L+U),\qquad \mathbf{g}=D^{-1}\mathbf{b}.
\]
矩阵 $B_J$ 称为 Jacobi 迭代矩阵. 给定初始向量 $\mathbf{x}_0=\bigl[x_1^{(0)},x_2^{(0)},\ldots,x_n^{(0)}\bigr]^T$ 后, 相应的迭代
\begin{equation}
  \mathbf{x}_k=B_J\mathbf{x}_{k-1}+\mathbf{g},\qquad k=1,2,\ldots,
  \label{eq:5-1}
\end{equation}
称为\newterm{jacobimethod}{Jacobi 方法}{Jacobi method}.

\subsection{Gauss-Seidel 方法}
\label{subsec:5-1-2}

在 Jacobi 方法中, 为计算向量
\[
  \mathbf{x}_{k+1}=\Bigl[x_1^{(k+1)},x_2^{(k+1)},\ldots,x_n^{(k+1)}\Bigr]^T
\]
的各分量, 只用到向量 $\mathbf{x}_k$ 的各分量. 然而注意, 为计算 $x_i^{(k+1)}$, 我们本可以使用已经得到的 $x_1^{(k+1)},x_2^{(k+1)},\ldots,x_{i-1}^{(k+1)}$. 因此, Jacobi 方法的一个自然改进是把 Jacobi 迭代 \eqref{eq:5-1} 改写为
\begin{equation}
  \mathbf{x}_k=(D-L)^{-1}U\mathbf{x}_{k-1}+(D-L)^{-1}\mathbf{b},\qquad k=1,2,\ldots
  \label{eq:5-2}
\end{equation}
其思想是: 每算出一个新分量, 就立即在下一个分量的计算中使用它. 迭代 \eqref{eq:5-2} 称为\newterm{gaussseidelmethod}{Gauss-Seidel 方法}{Gauss-Seidel method}.

注意矩阵 $D-L$ 是以 $a_{11},\ldots,a_{nn}$ 为对角元的下三角矩阵. 由于假定这些元素非零, 故矩阵 $D-L$ 非奇异. 矩阵
\[
  B_{GS}=(D-L)^{-1}U
\]
称为 Gauss-Seidel 迭代矩阵.

\section{收敛性分析}
\label{sec:5-2}

\subsection{收敛定理}
\label{subsec:5-2-1}

通常很难给出一个好的初始近似 $\mathbf{x}_0$. 因此, 最好能有条件保证 Jacobi 方法与 Gauss-Seidel 方法对任意选定的初始近似都收敛.

Jacobi 迭代与 Gauss-Seidel 迭代都可用下式表示
\begin{equation}
  \mathbf{x}_{k+1}=B\mathbf{x}_k+\mathbf{g},\qquad k=0,1,\ldots
  \label{eq:5-3}
\end{equation}
对 Jacobi 迭代, 有
\[
  B_J=D^{-1}(L+U),\qquad \mathbf{g}=D^{-1}\mathbf{b};
\]
而对 Gauss-Seidel 迭代, 有
\[
  B_{GS}=(D-L)^{-1}U,\qquad \mathbf{g}=(D-L)^{-1}\mathbf{b}.
\]
迭代 \eqref{eq:5-3} 称为\newterm{linearstationary}{线性定常迭代}{linear stationary iteration}, 其中 $B\in\R^{n\times n}$ 称为\newterm{iterationmatrix}{迭代矩阵}{iteration matrix}, $\mathbf{g}\in\R^n$ 称为常向量, $\mathbf{x}_0\in\R^n$ 称为初始向量. 下面给出一个收敛性定理.

\begin{theorem}
\label{thm:5-1}
迭代 \eqref{eq:5-3} 对任意初始猜测 $\mathbf{x}_0$ 收敛, 当且仅当 $k\to\infty$ 时 $B^k\to0$.
\end{theorem}

\begin{proof}
由 $\mathbf{x}=B\mathbf{x}+\mathbf{g}$ 与 $\mathbf{x}_{k+1}=B\mathbf{x}_k+\mathbf{g}$, 有
\begin{equation}
  \mathbf{x}-\mathbf{x}_{k+1}=B(\mathbf{x}-\mathbf{x}_k).
  \label{eq:5-4}
\end{equation}
由于它对 $k$ 的任意值都成立, 可以写出
\begin{equation}
  \mathbf{x}-\mathbf{x}_k=B(\mathbf{x}-\mathbf{x}_{k-1}).
  \label{eq:5-5}
\end{equation}
把 \eqref{eq:5-5} 代入 \eqref{eq:5-4}, 得
\[
  \mathbf{x}-\mathbf{x}_{k+1}=B^2(\mathbf{x}-\mathbf{x}_{k-1}).
\]
如此继续 $k$ 次, 可写出
\[
  \mathbf{x}-\mathbf{x}_{k+1}=B^{k+1}(\mathbf{x}-\mathbf{x}_0).
\]
这表明, 对任意选定的 $\mathbf{x}_0$, $\{\mathbf{x}_k\}$ 收敛到解 $\mathbf{x}$ 当且仅当 $k\to\infty$ 时 $B^k\to0$.
\end{proof}

回忆定理 \ref{thm:3-5}: $k\to\infty$ 时 $B^k\to0$ 当且仅当谱半径 $\rho(B)<1$. 由于 $|\lambda_i|\leqslant\norm{B}$, 判断 $\rho(B)<1$ 的一个好办法是用行和范数或列和范数计算 $\norm{B}$, 看是否有 $\norm{B}<1$. 注意其逆不真. 把定理 \ref{thm:5-1} 的结论与上述观察结合起来, 便得下述定理.

\begin{theorem}
\label{thm:5-2}
迭代 \eqref{eq:5-3} 对任意选定的 $\mathbf{x}_0$ 收敛, 当且仅当 $\rho(B)<1$. 此外, 若对某个矩阵范数有 $\norm{B}<1$, 则迭代 \eqref{eq:5-3} 收敛.
\end{theorem}

考虑下面的例子:
\[
  A_1=
  \begin{bmatrix}
    1 & 2 & -2\\
    1 & 1 & 1\\
    2 & 2 & 1
  \end{bmatrix},
  \qquad
  A_2=
  \begin{bmatrix}
    2 & -1 & 1\\
    1 & 1 & 1\\
    1 & 1 & -2
  \end{bmatrix}.
\]
容易验证, 对 $A_1$, 即使 Gauss-Seidel 方法不收敛, Jacobi 方法仍收敛; 对 $A_2$, Jacobi 方法发散而 Gauss-Seidel 方法收敛.

\subsection{收敛的充分条件}
\label{subsec:5-2-2}

现在应用定理 \ref{thm:5-2} 给出一系列准则, 它们保证 Jacobi 方法和 (或)Gauss-Seidel 方法对任意选定的初始近似 $\mathbf{x}_0$ 收敛.

\begin{definition}
\label{def:5-1}
设 $A=[a_{ij}]\in\R^{n\times n}$. 若
\[
  |a_{ii}|>\sum_{\substack{j=1\\ j\neq i}}^{n}|a_{ij}|,
  \qquad i=1,2,\ldots,n,
\]
则称矩阵 $A$ 为严格对角占优的. 若
\[
  |a_{ii}|\geqslant\sum_{\substack{j=1\\ j\neq i}}^{n}|a_{ij}|,
  \qquad i=1,2,\ldots,n,
\]
且其中至少有一个严格不等式成立, 则称矩阵 $A$ 为\newterm{weaklydiagdominant}{弱对角占优的}{weakly diagonally dominant}.
\end{definition}

\begin{theorem}
\label{thm:5-3}
若 $A$ 严格对角占优, 则 Jacobi 方法对任意初始近似 $\mathbf{x}_0$ 收敛.
\end{theorem}

\begin{proof}
由于 $A=[a_{ij}]$ 严格对角占优, 有
\[
  |a_{ii}|>\sum_{\substack{j=1\\ j\neq i}}^{n}|a_{ij}|,
  \qquad i=1,2,\ldots,n.
\]
回忆 Jacobi 迭代矩阵
\[
  B_J=D^{-1}(L+U)
\]
由下式给出:
\[
  B_J=
  \begin{bmatrix}
    0 & -\dfrac{a_{12}}{a_{11}} & \cdots & \cdots & -\dfrac{a_{1n}}{a_{11}}\\[8pt]
    -\dfrac{a_{21}}{a_{22}} & 0 & -\dfrac{a_{23}}{a_{22}} & \cdots & -\dfrac{a_{2n}}{a_{22}}\\[8pt]
    \vdots & \ddots & \ddots & \ddots & \vdots\\[2pt]
    \vdots & & \ddots & \ddots & -\dfrac{a_{n-1,n}}{a_{n-1,n-1}}\\[8pt]
    -\dfrac{a_{n1}}{a_{nn}} & \cdots & \cdots & \dfrac{a_{n,n-1}}{a_{nn}} & 0
  \end{bmatrix}.
\]
我们知道每一行的绝对行和都小于 $1$, 即
\[
  \norm{B_J}_\infty<1.
\]
于是由定理 \ref{thm:5-2} 知, Jacobi 方法收敛.
\end{proof}

\begin{theorem}
\label{thm:5-4}
若 $A=[a_{ij}]\in\R^{n\times n}$ 严格对角占优, 则 Gauss-Seidel 方法对任意初始近似 $\mathbf{x}_0$ 收敛.
\end{theorem}

\begin{proof}
Gauss-Seidel 迭代矩阵由
\[
  B_{GS}=(D-L)^{-1}U
\]
给出. 设 $\lambda$ 是该矩阵的特征值, $\mathbf{x}=[x_1,x_2,\ldots,x_n]^T$ 是相应的特征向量, 其模最大的分量为 $1$. 则由
\[
  B_{GS}\mathbf{x}=\lambda\mathbf{x}
\]
得
\[
  U\mathbf{x}=\lambda(D-L)\mathbf{x},
\]
即
\[
  -\sum_{j=i+1}^{n}a_{ij}x_j=\lambda\sum_{j=1}^{i}a_{ij}x_j,
  \qquad 1\leqslant i\leqslant n,
\]
它可改写为
\begin{equation}
  \lambda a_{ii}x_i=-\lambda\sum_{j=1}^{i-1}a_{ij}x_j-\sum_{j=i+1}^{n}a_{ij}x_j,
  \qquad 1\leqslant i\leqslant n.
  \label{eq:5-6}
\end{equation}
设 $x_k$ 是向量 $\mathbf{x}$ 中模最大的分量 (其模为 $1$). 则由 \eqref{eq:5-6} 得
\[
  |\lambda|\cdot|a_{kk}|\leqslant|\lambda|\sum_{j=1}^{k-1}|a_{kj}|+\sum_{j=k+1}^{n}|a_{kj}|,
\]
即
\[
  |\lambda|\Bigl(|a_{kk}|-\sum_{j=1}^{k-1}|a_{kj}|\Bigr)\leqslant\sum_{j=k+1}^{n}|a_{kj}|,
\]
或
\begin{equation}
  |\lambda|\leqslant
  \frac{\displaystyle\sum_{j=k+1}^{n}|a_{kj}|}
       {\displaystyle|a_{kk}|-\sum_{j=1}^{k-1}|a_{kj}|}.
  \label{eq:5-7}
\end{equation}
由于 $A$ 严格对角占优, 有
\[
  |a_{kk}|-\sum_{j=1}^{k-1}|a_{kj}|>\sum_{j=k+1}^{n}|a_{kj}|.
\]
于是由 \eqref{eq:5-7} 断定 $|\lambda|<1$, 即 $\rho(B_{GS})<1$. 由定理 \ref{thm:5-2} 知, Gauss-Seidel 方法收敛.
\end{proof}

下面讨论对称正定矩阵的 Jacobi 方法与 Gauss-Seidel 方法的收敛性.

\begin{theorem}
\label{thm:5-5}
设 $A=[a_{ij}]\in\R^{n\times n}$ 对称, 且对角元 $a_{ii}>0$ ($i=1,2,\ldots,n$). 则 Jacobi 方法收敛当且仅当 $A$ 与 $2D-A$ 都正定.
\end{theorem}

\begin{proof}
由于
\[
  B_J=D^{-1}(L+U)=D^{-1}(D-A)=I-D^{-1}A,
\]
且
\[
  D=\diag(a_{11},a_{22},\ldots,a_{nn})
\]
中 $a_{ii}>0$ ($i=1,2,\ldots,n$), 故
\[
  B_J=I-D^{-1}A=D^{-1/2}\bigl(I-D^{-1/2}AD^{-1/2}\bigr)D^{1/2}.
\]
容易看出 $I-D^{-1/2}AD^{-1/2}$ 是对称的, 且与 $B_J$ 相似. 于是 $B_J$ 的特征值均为实数.

现在先假设 Jacobi 方法收敛, 则由定理 \ref{thm:5-2} 知 $\rho(B_J)<1$. 矩阵
\[
  I-D^{-1/2}AD^{-1/2}
\]
的特征值的绝对值小于 $1$, 即 $D^{-1/2}AD^{-1/2}$ 的特征值落在 $(0,2)$ 内. 因此 $A$ 正定. 另一方面, $2I-D^{-1/2}AD^{-1/2}$ 的特征值为正, 故矩阵
\[
  2I-D^{-1/2}AD^{-1/2}
\]
正定. 由于
\[
  D^{-1/2}(2D-A)D^{-1/2}=2I-D^{-1/2}AD^{-1/2},
\]
可知 $2D-A$ 也正定.

反之, 由于
\[
  D^{1/2}(I-B_J)D^{-1/2}=D^{-1/2}AD^{-1/2},
\]
且 $A$ 正定, 推知 $I-B_J$ 的特征值为正, 即 $B_J$ 的特征值小于 $1$. 又因 $2D-A$ 正定, 且
\[
  D^{-1/2}(2D-A)D^{-1/2}=D^{1/2}\bigl(I+B_J\bigr)D^{-1/2},
\]
可推知 $I+B_J$ 的特征值为正, 即 $B_J$ 的特征值大于 $-1$. 于是 $\rho(B_J)<1$. 再次由定理 \ref{thm:5-2} 知, Jacobi 方法收敛.
\end{proof}

\begin{theorem}
\label{thm:5-6}
设 $A$ 为对称正定矩阵. 则 Gauss-Seidel 方法对任意初始近似 $\mathbf{x}_0$ 收敛.
\end{theorem}

\begin{proof}
设 $\lambda$ 是迭代矩阵 $B_{GS}$ 的特征值, $\mathbf{u}$ 是相应的特征向量. 则
\[
  (D-L)^{-1}U\mathbf{u}=\lambda\mathbf{u}.
\]
由于 $A$ 对称正定, 有 $U=L^T$, 且
\[
  \lambda(D-L)\mathbf{u}=L^T\mathbf{u}.
\]
因此
\[
  \lambda\mathbf{u}^*(D-L)\mathbf{u}=\mathbf{u}^*L^T\mathbf{u}.
\]
令
\[
  \mathbf{u}^*D\mathbf{u}=\delta,\qquad \mathbf{u}^*L\mathbf{u}=\alpha+\ii\beta.
\]
于是有
\[
  \mathbf{u}^*L^T\mathbf{u}=(L\mathbf{u})^*\mathbf{u}=\overline{\mathbf{u}^*L\mathbf{u}}=\alpha-\ii\beta.
\]
因此
\[
  \lambda\bigl[\delta-(\alpha+\ii\beta)\bigr]=\alpha-\ii\beta.
\]
两边取模, 得
\[
  |\lambda|^2=\frac{\alpha^2+\beta^2}{(\delta-\alpha)^2+\beta^2}.
\]
另一方面,
\[
  0<\mathbf{u}^*A\mathbf{u}=\mathbf{u}^*(D-L-L^T)\mathbf{u}=\delta-2\alpha.
\]
因此
\[
  \begin{aligned}
    (\delta-\alpha)^2+\beta^2
    &=\delta^2+\alpha^2+\beta^2-2\delta\alpha\\
    &=\delta(\delta-2\alpha)+\alpha^2+\beta^2\\
    &>\alpha^2+\beta^2.
  \end{aligned}
\]
于是得到 $|\lambda|<1$. 由定理 \ref{thm:5-2} 知, Gauss-Seidel 方法收敛.
\end{proof}

\begin{definition}
\label{def:5-2}
若不存在置换矩阵 $P$ 使得
\[
  PAP^T=
  \begin{bmatrix}
    A_{11} & A_{12}\\
    \mathbf{0} & A_{22}
  \end{bmatrix},
\]
其中 $A_{11}$ 与 $A_{22}$ 为方阵, 则称矩阵 $A$ 是\newterm{irreducible}{不可约的}{irreducible}.
\end{definition}

我们有下述引理 \cite{src:85}.

\begin{lemma}
\label{lem:5-1}
若矩阵 $A$ 严格对角占优; 或者 $A$ 不可约且弱对角占优, 则 $A$ 非奇异.
\end{lemma}

此外,

\begin{theorem}
\label{thm:5-7}
我们有
\begin{romanlist}
  \item 若 $A$ 严格对角占优, 则 Jacobi 方法与 Gauss-Seidel 方法都收敛. 事实上,
  \[
    \norm{B_{GS}}_\infty\leqslant\norm{B_J}_\infty<1.
  \]
  \item 若 $A$ 不可约且弱对角占优, 则 Jacobi 方法与 Gauss-Seidel 方法都收敛. 此外,
  \[
    \rho(B_{GS})<\rho(B_J)<1.
  \]
\end{romanlist}
\end{theorem}

本定理的证明参见 \cite{src:30,src:85}.

\section{收敛速度}
\label{sec:5-3}

我们考虑线性定常迭代法
\[
  \mathbf{x}_{k+1}=M\mathbf{x}_k+\mathbf{g},\qquad k=0,1,\ldots,
\]
其中 $M$ 是 $n\times n$ 矩阵. 若 $I-M$ 非奇异, 则
\[
  (I-M)\mathbf{x}_*=\mathbf{g}
\]
有唯一解. 误差向量 $\mathbf{y}_k$ 定义为 $\mathbf{y}_k=\mathbf{x}_k-\mathbf{x}_*$, $k=0,1,\ldots$, 于是
\[
  \mathbf{y}_k=M\mathbf{y}_{k-1}=\cdots=M^k\mathbf{y}_0.
\]
利用矩阵范数与向量范数, 有
\[
  \norm{\mathbf{y}_k}\leqslant\norm{M^k}\,\norm{\mathbf{y}_0},
\]
且对每个 $k$, 都可能对某个向量 $\mathbf{y}_0$ 取等号. 因此, 若 $\mathbf{y}_0$ 不是零向量, 则 $\norm{M^k}$ 给出了比值 $\norm{\mathbf{y}_k}/\norm{\mathbf{y}_0}$ 的一个精确上界估计. 由于实际问题中初始向量 $\mathbf{y}_0$ 未知, 故 $\norm{M^k}$ 可作为比较不同迭代法的一种度量.

\begin{definition}
\label{def:5-3}
设 $M$ 为 $n\times n$ 迭代矩阵. 若对某个正整数 $k$ 有 $\norm{M^k}<1$, 则量
\[
  R_k(M)\equiv-\ln\bigl[(\norm{M^k})^{1/k}\bigr]=-\frac{\ln\norm{M^k}}{k}
\]
称为 $M$ 迭代 $k$ 次的\newterm{averagerate}{平均收敛速度}{average rate of convergence}.
\end{definition}

从实际计算的角度看, 平均收敛速度 $R_k(M)$ 的意义如下. 显然, 量
\[
  \sigma=\left(\frac{\norm{\mathbf{y}_k}}{\norm{\mathbf{y}_0}}\right)^{1/k}
\]
是误差范数每次迭代的平均缩减因子. 若 $\norm{M^k}<1$, 则由定义 \ref{def:5-3} 有
\[
  \sigma\leqslant(\norm{M^k})^{1/k}=e^{-R_k(M)},
\]
其中 $e$ 是自然对数的底.

若 $M$ 是对称的 (或 Hermitian 的, 或正规的, 即 $MM^*=M^*M$), 则利用迭代矩阵 $M$ 的谱半径, 有
\[
  \norm{M^k}_2=[\rho(M)]^k,
\]
从而
\[
  R_k(M)=-\ln\rho(M),
\]
它与 $k$ 无关.

下面考虑\newterm{asymptoticrate}{渐近收敛速度}{asymptotic convergence rate}
\[
  R_\infty(M)\equiv\lim_{k\to\infty}R_k(M).
\]

\begin{theorem}
\label{thm:5-8}
有
\[
  R_\infty(M)=-\ln\rho(M).
\]
\end{theorem}

\begin{proof}
只需证明
\[
  \lim_{k\to\infty}\norm{M^k}^{1/k}=\rho(M).
\]
由于
\[
  [\rho(M)]^k=\rho(M^k)\leqslant\norm{M^k},
\]
有
\[
  \rho(M)\leqslant\norm{M^k}^{1/k}.
\]
另一方面, 对任意 $\epsilon>0$, 考虑矩阵
\[
  B_\epsilon=\frac{1}{\rho(M)+\epsilon}M.
\]
显然 $\rho(B_\epsilon)<1$, 从而 $\lim_{k\to\infty}B_\epsilon^k=0$. 因此存在自然数 $K$, 当 $k\geqslant K$ 时有
\[
  \norm{B_\epsilon^k}\leqslant1,
\]
即
\[
  \norm{M^k}\leqslant[\rho(M)+\epsilon]^k.
\]
于是
\[
  \rho(M)\leqslant\norm{M^k}^{1/k}\leqslant\rho(M)+\epsilon,
\]
这意味着
\[
  \lim_{k\to\infty}\norm{M^k}^{1/k}=\rho(M).\qedhere
\]
\end{proof}

\section{SOR 方法}
\label{sec:5-4}

当 $\rho(B_{GS})$ 接近 $1$ 时, Gauss-Seidel 方法非常慢. 然而在某些情形下, 引入一个称为\newterm{relaxationparameter}{松弛因子}{relaxation parameter}的参数 $\omega$ 可以提高 Gauss-Seidel 迭代的收敛速度. 这种方法称为\newterm{successiveoverrelaxation}{逐次超松弛方法}{successive overrelaxation (SOR) method}.

\subsection{迭代形式}
\label{subsec:5-4-1}

设 $A=[a_{ij}]\in\R^{n\times n}$, 且 $A=D-L-U$ 如 \ref{subsec:5-1-1} 小节所定义. 再次考虑 $A\mathbf{x}=\mathbf{b}$ 的求解. SOR 方法的动机是对 $x_i^{(k)}$ 与 $x_i^{(k+1)}$ 取适当的加权平均来改进 Gauss-Seidel 迭代, 即
\begin{equation}
  x_i^{(k+1)}=(1-\omega)x_i^{(k)}
  +\omega\Biggl(\sum_{j=1}^{i-1}c_{ij}x_j^{(k+1)}
  +\sum_{j=i+1}^{n}c_{ij}x_j^{(k)}+g_i\Biggr).
  \label{eq:5-8}
\end{equation}
这里 $D^{-1}(L+U)=[c_{ij}]$, $\mathbf{x}_k=\bigl[x_1^{(k)},x_2^{(k)},\ldots,x_n^{(k)}\bigr]^T$, 且 $\mathbf{g}=D^{-1}\mathbf{b}=[g_1,g_2,\ldots,g_n]^T$. 写成矩阵形式, 有
\[
  \mathbf{x}_{k+1}=L_\omega\mathbf{x}_k+\omega(D-\omega L)^{-1}\mathbf{b},
\]
其中
\[
  L_\omega=(D-\omega L)^{-1}\bigl[(1-\omega)D+\omega U\bigr]
\]
称为 SOR 方法的迭代矩阵, $\omega$ 为松弛因子.

按 $\omega$ 的取值有三种情形:
\begin{parenlist}
  \item 若 $\omega=1$, 则 \eqref{eq:5-8} 等价于 Gauss-Seidel 方法;
  \item 若 $\omega<1$, 则 \eqref{eq:5-8} 称为\newterm{underrelaxation}{低松弛}{underrelaxation};
  \item 若 $\omega>1$, 则 \eqref{eq:5-8} 称为\newterm{overrelaxation}{超松弛}{overrelaxation}.
\end{parenlist}

\subsection{收敛准则}
\label{subsec:5-4-2}

自然会问: 对 $\omega$ 的哪些取值范围, SOR 迭代收敛? 为此, 我们先证明下述重要结果 \cite{src:30,src:85}.

\begin{theorem}
\label{thm:5-9}
若 $\omega$ 落在区间 $(0,2)$ 之外, 则 SOR 迭代对任何初始近似都不收敛.
\end{theorem}

\begin{proof}
回忆 SOR 迭代矩阵 $L_\omega$ 由
\[
  L_\omega=(D-\omega L)^{-1}\bigl[(1-\omega)D+\omega U\bigr]
\]
给出, 其中 $A=[a_{ij}]=D-L-U$.

矩阵 $(D-\omega L)^{-1}$ 是以 $1/a_{ii}$ ($i=1,2,\ldots,n$) 为对角元的下三角矩阵, 而矩阵
\[
  (1-\omega)D+\omega U
\]
是以 $(1-\omega)a_{ii}$ ($i=1,2,\ldots,n$) 为对角元的上三角矩阵. 因此
\[
  \det(L_\omega)=(1-\omega)^n.
\]
由于矩阵的行列式等于其特征值之积, 我们断定
\[
  \rho(L_\omega)\geqslant|1-\omega|,
\]
其中 $\rho(L_\omega)$ 是 $L_\omega$ 的谱半径. 由定理 \ref{thm:5-2} 知, 迭代矩阵的谱半径应小于 $1$ 才能收敛. 于是断定: SOR 方法的收敛要求 $0<\omega<2$.
\end{proof}

\begin{theorem}
\label{thm:5-10}
若 $A$ 对称正定, 则 $\rho(L_\omega)<1$ 对一切 $0<\omega<2$ 成立.
\end{theorem}

\begin{proof}
设 $\lambda$ 是 SOR 迭代矩阵
\[
  L_\omega=(D-\omega L)^{-1}\bigl[(1-\omega)D+\omega L^T\bigr]
\]
的任一特征值, $\mathbf{x}$ 是相应的特征向量. 有
\[
  \bigl[(1-\omega)D+\omega L^T\bigr]\mathbf{x}=\lambda(D-\omega L)\mathbf{x},
\]
或
\[
  \mathbf{x}^*\bigl[(1-\omega)D+\omega L^T\bigr]\mathbf{x}
  =\lambda\mathbf{x}^*(D-\omega L)\mathbf{x}.
\]
令
\[
  \mathbf{x}^*D\mathbf{x}=\delta,\qquad \mathbf{x}^*L\mathbf{x}=\alpha+\ii\beta.
\]
因此 $\mathbf{x}^*L^T\mathbf{x}=\alpha-\ii\beta$, 从而
\[
  (1-\omega)\delta+\omega(\alpha-\ii\beta)
  =\lambda\bigl[\delta-\omega(\alpha+\ii\beta)\bigr].
\]
两边取模, 得
\begin{equation}
  |\lambda|^2=
  \frac{\bigl[(1-\omega)\delta+\omega\alpha\bigr]^2+\omega^2\beta^2}
       {(\delta-\omega\alpha)^2+\omega^2\beta^2}.
  \label{eq:5-9}
\end{equation}
注意
\[
  \begin{aligned}
    &\bigl[(1-\omega)\delta+\omega\alpha\bigr]^2+\omega^2\beta^2
     -(\delta-\omega\alpha)^2-\omega^2\beta^2\\
    &\quad=\bigl[\delta-\omega(\delta-\alpha)\bigr]^2-(\delta-\omega\alpha)^2\\
    &\quad=\omega\delta(\delta-2\alpha)(\omega-2).
  \end{aligned}
\]
由于 $A$ 对称正定, 有
\[
  \delta>0,\qquad \delta-2\alpha>0.
\]
因此, 若 $0<\omega<2$, 则
\begin{equation}
  \bigl[(1-\omega)\delta+\omega\alpha\bigr]^2+\omega^2\beta^2
  <(\delta-\omega\alpha)^2+\omega^2\beta^2.
  \label{eq:5-10}
\end{equation}
于是, 对 $0<\omega<2$, 由 \eqref{eq:5-9} 与 \eqref{eq:5-10} 得
\[
  |\lambda|^2<1,
\]
即 SOR 方法收敛.
\end{proof}

\subsection{SOR 迭代中的最优 $\omega$}
\label{subsec:5-4-3}

为了进一步比较 Jacobi, Gauss-Seidel 与 SOR 方法, 我们对矩阵再施加一个条件. 这一条件使我们能用 $\rho(B_J)$ 显式地算出 $\rho(B_{GS})$ 与 $\rho(L_\omega)$.

\begin{definition}
\label{def:5-4}
若存在置换 $P$ 使得
\[
  PMP^T=
  \begin{bmatrix}
    D_{11} & M_{12}\\
    M_{21} & D_{22}
  \end{bmatrix},
\]
其中 $D_{11}$ 与 $D_{22}$ 为对角矩阵, 则称矩阵 $M$ 具有\newterm{propertya}{性质 A}{property A}.
\end{definition}

若 $M$ 具有性质 A, 则可写
\[
  PMP^T=\widehat{D}-\widehat{L}-\widehat{U},
\]
其中
\[
  \widehat{D}=
  \begin{bmatrix}
    D_{11} & \mathbf{0}\\
    \mathbf{0} & D_{22}
  \end{bmatrix},
  \qquad
  \widehat{L}=
  \begin{bmatrix}
    \mathbf{0} & \mathbf{0}\\
    -M_{21} & \mathbf{0}
  \end{bmatrix},
  \qquad
  \widehat{U}=
  \begin{bmatrix}
    \mathbf{0} & -M_{12}\\
    \mathbf{0} & \mathbf{0}
  \end{bmatrix}.
\]
令
\[
  \widehat{B_J}(\alpha)\equiv\alpha\widehat{D}^{-1}\widehat{L}
  +\frac{1}{\alpha}\widehat{D}^{-1}\widehat{U}.
\]
我们有

\begin{theorem}
\label{thm:5-11}
$\widehat{B_J}(\alpha)$ 的特征值与 $\alpha$ 无关.
\end{theorem}

\begin{proof}
只需注意矩阵
\[
  \widehat{B_J}(\alpha)=-
  \begin{bmatrix}
    \mathbf{0} & \dfrac{1}{\alpha}D_{11}^{-1}M_{12}\\[8pt]
    \alpha D_{22}^{-1}M_{21} & \mathbf{0}
  \end{bmatrix}
\]
与矩阵
\[
  \begin{bmatrix}
    I & \mathbf{0}\\
    \mathbf{0} & \alpha I
  \end{bmatrix}^{-1}
  \widehat{B_J}(\alpha)
  \begin{bmatrix}
    I & \mathbf{0}\\
    \mathbf{0} & \alpha I
  \end{bmatrix}
  =-
  \begin{bmatrix}
    \mathbf{0} & D_{11}^{-1}M_{12}\\
    D_{22}^{-1}M_{21} & \mathbf{0}
  \end{bmatrix}
  =\widehat{B_J}(1).\qedhere
\]
\end{proof}

\begin{definition}
\label{def:5-5}
设 $M=D-L-U$, 且
\[
  B_J(\alpha)=\alpha D^{-1}L+\frac{1}{\alpha}D^{-1}U.
\]
若 $B_J(\alpha)$ 的特征值与 $\alpha$ 无关, 则称 $M$ 具有\newterm{consistentordering}{相容次序}{consistent ordering}(或称 $M$ 是相容次序的).
\end{definition}

注意 $B_J(1)=B_J$ 是 Jacobi 迭代矩阵. 由定理 \ref{thm:5-11} 知, 若 $M$ 具有性质 A, 则 $PMP^T$ 是相容次序的, 其中 $P$ 是使
\[
  PMP^T=
  \begin{bmatrix}
    D_{11} & M_{12}\\
    M_{21} & D_{22}
  \end{bmatrix}
\]
的置换矩阵, 而 $D_{11}$ 与 $D_{22}$ 为对角矩阵. 相容次序并不蕴涵性质 A.

\begin{example}
\label{ex:5-1}
当 $D_i$ 为对角矩阵时, 任何\newterm{blocktridiagonal}{块三对角矩阵}{block tridiagonal matrix}
\[
  \begin{bmatrix}
    D_1 & A_1 & & \\
    B_1 & \ddots & \ddots & \\
      & \ddots & \ddots & A_{n-1}\\
      & & B_{n-1} & D_n
  \end{bmatrix}
\]
都是相容次序的.
\end{example}

下述定理给出 $B_J$ 的特征值与 $L_\omega$ 的特征值之间的关系 \cite{src:30,src:85}.

\begin{theorem}
\label{thm:5-12}
若 $A$ 是相容次序的, 且 $\omega\neq0$, 则
\begin{romanlist}
  \item $B_J$ 的特征值以 $\pm$ 对的形式出现.
  \item 若 $\mu$ 是 $B_J$ 的特征值, 且 $\lambda$ 满足
  \begin{equation}
    (\lambda+\omega-1)^2=\lambda\omega^2\mu^2,
    \label{eq:5-11}
  \end{equation}
  则 $\lambda$ 是 $L_\omega$ 的特征值.
  \item 若 $\lambda\neq0$ 是 $L_\omega$ 的特征值, 且 $\mu$ 满足 \eqref{eq:5-11}, 则 $\mu$ 是 $B_J$ 的特征值.
\end{romanlist}
\end{theorem}

\begin{corollary}
\label{cor:5-1}
若 $A$ 是相容次序的, 则
\[
  \rho(B_{GS})=[\rho(B_J)]^2.
\]
这意味着 Gauss-Seidel 方法的收敛速度是 Jacobi 方法的两倍.
\end{corollary}

\begin{proof}
SOR 方法中取 $\omega=1$ 等价于 Gauss-Seidel 方法. 因此由 \eqref{eq:5-11} 有 $\lambda^2=\lambda\mu^2$, 从而
\[
  \lambda=\mu^2.\qedhere
\]
\end{proof}

为了从超松弛中获得最大收益, 我们希望找一个最优的 $\omega$, 记为 $\omega_{opt}$, 使 $\rho(L_\omega)$ 达到最小. 我们有下述定理 \cite{src:30,src:85}.

\begin{theorem}
\label{thm:5-13}
设 $A$ 是相容次序的, 且 $B_J$ 的特征值为实数, $\mu=\rho(B_J)<1$. 则
\[
  \omega_{opt}=\frac{2}{1+\sqrt{1-\mu^2}},
\]
\[
  \rho(L_{\omega_{opt}})=\frac{\mu^2}{\bigl(1+\sqrt{1-\mu^2}\bigr)^2},
\]
以及
\[
  \rho(L_\omega)=
  \begin{cases}
    \omega-1, & \omega_{opt}\leqslant\omega<2,\\[6pt]
    1-\omega+\dfrac{1}{2}\omega^2\mu^2
    +\omega\mu\sqrt{1-\omega+\dfrac{1}{4}\omega^2\mu^2},
    & 0<\omega\leqslant\omega_{opt}.
  \end{cases}
\]
\end{theorem}

关于经典迭代法更多的收敛性分析, 可参见 \cite{src:14,src:74}.

\begin{exercises}
  \item 判断下列例子中 Jacobi 方法与 Gauss-Seidel 方法的收敛性:
  \[
    A_1=
    \begin{bmatrix}
      1 & 2 & -2\\
      1 & 1 & 1\\
      2 & 2 & 1
    \end{bmatrix},
    \qquad
    A_2=
    \begin{bmatrix}
      2 & -1 & 1\\
      1 & 1 & 1\\
      1 & 1 & -2
    \end{bmatrix}.
  \]

  \item 证明对 $2\times2$ 对称正定方程组, Jacobi 方法收敛.

  \item 证明: 若 $A=M-N$ 奇异, 则即使 $M$ 非奇异, 也绝不可能有 $\rho(M^{-1}N)<1$.

  \item 考虑 $A\mathbf{x}=\mathbf{b}$, 其中
  \[
    A=
    \begin{bmatrix}
      1 & 0 & \alpha\\
      0 & 1 & 0\\
      \alpha & 0 & 1
    \end{bmatrix}.
  \]
  \subquest 对哪些 $\alpha$, $A$ 正定?
  \subquest 对哪些 $\alpha$, Jacobi 方法收敛?
  \subquest 对哪些 $\alpha$, Gauss-Seidel 方法收敛?

  \item 设 $A\in\R^{n\times n}$ 非奇异. 证明存在置换矩阵 $P$, 使得 $PA$ 的对角元均非零.

  \item 证明
  \[
    A=
    \begin{bmatrix}
      4 & -1 & -1 & 0\\
      -1 & 4 & 0 & -1\\
      -1 & 0 & 4 & -1\\
      0 & -1 & -1 & 4
    \end{bmatrix}
  \]
  是相容次序的.

  \item 设 $A\in\R^{n\times n}$. 则 $\rho(A)<1$ 当且仅当 $I-A$ 非奇异, 且 $(I-A)^{-1}(I+A)$ 的每个特征值都有正实部.

  \item 设 $B\in\R^{n\times n}$ 满足 $\rho(B)=0$. 证明: 对任意 $\mathbf{g},\mathbf{x}_0\in\R^n$, 迭代公式
  \[
    \mathbf{x}_{k+1}=B\mathbf{x}_k+\mathbf{g},\qquad k=0,1,\ldots
  \]
  至多经过 $n$ 次迭代就收敛到 $\mathbf{x}=B\mathbf{x}+\mathbf{g}$ 的精确解.

  \item 若 $A=[a_{ij}]\in\R^{n\times n}$ 不可约, $a_{ij}\geqslant0$, 而 $B=[b_{ij}]\in\R^{n\times n}$, $b_{ij}\geqslant0$. 证明 $A+B$ 不可约.

  \item 设
  \[
    A=
    \begin{bmatrix}
      0 & a & 0\\
      0 & 0 & b\\
      c & 0 & 0
    \end{bmatrix}.
  \]
  $A^k$ 是否不可约 ($k=1,2,3$)?

  \item 设
  \[
    A=
    \begin{bmatrix}
      2 & -1\\
      1 & 0
    \end{bmatrix}.
  \]
  证明
  \subquest
  \[
    A^k=k
    \begin{bmatrix}
      1+1/k & -1\\
      1 & 1/k-1
    \end{bmatrix},
    \qquad k=1,2,\ldots
  \]
  \subquest
  \[
    \lim_{k\to\infty}\frac{\norm{A^k}_p}{k}=2,
    \qquad p=1,2,\infty.
  \]
  \subquest $\rho(A^k)=1$.

  \item 设
  \[
    B_k=B_{k-1}+B_{k-1}(I-AB_{k-1}),\qquad k=1,2,\ldots
  \]
  证明: 若 $\norm{I-AB_0}=c<1$, 则
  \[
    \lim_{k\to\infty}B_k=A^{-1}
  \]
  且
  \[
    \norm{A^{-1}-B_k}\leqslant\frac{c^{2^k}}{1-c}\norm{B_0}.
  \]

  \item 证明定理 \ref{thm:5-12}.

  \item 证明定理 \ref{thm:5-13}.
\end{exercises}
